% =========================================================
% The Mediant Property, Farey Sequences, and Stern--Brocot Theory
% =========================================================

\subsection{The Mediant Property}
\label{sec:mediant}

\begin{tcolorbox}[
  colback=gray!6,
  colframe=gray!40,
  arc=2pt,
  left=8pt, right=8pt, top=6pt, bottom=6pt,
  title={\small\textbf{Toolkit: Mediants, Farey Sequences, and Stern--Brocot}},
  fonttitle=\small\bfseries
]
Once density is known abstractly, it is natural to ask for a \emph{constructive}
way to produce rationals between two given ones, and to understand the precise
arithmetic of the gaps between them.  The atomic notions developed here are:
\begin{itemize}
  \item the \textbf{mediant}, and its order-theoretic property;
  \item \textbf{Farey sequences}, the canonical enumeration of rationals by
        denominator bound;
  \item the \textbf{Farey Neighbor Theorem} and the exact gap formula;
  \item the \textbf{Mediant Gap Lemma}, quantifying what happens when the
        mediant is inserted;
  \item the \textbf{Stern--Brocot tree}, which enumerates all positive rationals
        without repetition.
\end{itemize}
\end{tcolorbox}

% ---------------------------------------------------------
\subsubsection{Mediants}
% ---------------------------------------------------------

\begin{tcolorbox}[colback=defbox, colframe=defborder, arc=2pt,
  left=6pt, right=6pt, top=4pt, bottom=4pt,
  title={\small\textbf{Definition (Mediant)}},
  fonttitle=\small\bfseries]
\begin{definition}[Mediant]
\label{def:mediant}
Let $a/b, c/d \in \mathbb{Q}$ with $b, d > 0$.
The \emph{mediant} of these two fractions is
\[
\frac{a+c}{b+d}.
\]
\end{definition}
\end{tcolorbox}

\begin{remark*}[Standard quantified statement]
\[
\forall a,c \in \mathbb{Z},\; \forall b,d \in \mathbb{Z}_{>0},\quad
\operatorname{Mediant}\!\left(\tfrac{a}{b},\tfrac{c}{d}\right)
= \frac{a+c}{b+d}.
\]
\end{remark*}

\begin{remark*}[Interpretation]
The mediant is formed by adding numerators and denominators separately.
This is not fraction addition — it is a distinct operation that always
produces a value strictly between its two inputs when they are in order.
Its strength is constructive: given any two rationals, the mediant produces
a specific third rational between them with full arithmetic control over
where it lands and what gap it leaves.
\end{remark*}

\begin{remark*}[Dependencies]
Order and arithmetic of $\mathbb{Q}$,
\hyperref[def:rational-operations]{rational operations}.
\end{remark*}

\begin{theorem}[Mediant inequality]
\label{thm:mediant-inequality}
If $a/b < c/d$ with $b, d > 0$, then
\[
\frac{a}{b} < \frac{a+c}{b+d} < \frac{c}{d}.
\]
\hyperref[prf:mediant-inequality]{Go to proof.}
\end{theorem}

\begin{remark*}[Standard quantified statement]
\[
\forall a,c \in \mathbb{Z},\; \forall b,d \in \mathbb{Z}_{>0},\quad
\frac{a}{b} < \frac{c}{d}
\;\Longrightarrow\;
\frac{a}{b} < \frac{a+c}{b+d} < \frac{c}{d}.
\]
\end{remark*}

\begin{remark*}[Interpretation]
The mediant always lands strictly inside the interval.  This makes it a
concrete construction proof of density: given any two rationals, produce
a third one between them without appeal to existence arguments.
\end{remark*}

\begin{remark*}[Dependencies]
\begin{itemize}
  \item \hyperref[def:mediant]{Mediant}
  \item \hyperref[thm:q-is-ordered-field]{$\mathbb{Q}$ Is an Ordered Field}
\end{itemize}
\end{remark*}

\begin{remark}[Local versus global density]
\label{rem:mediant-vs-arch}
The mediant property is \emph{local}: it refines a given interval by
inserting a specific rational between two given rationals.
The Archimedean property is \emph{global}: it says the natural numbers
eventually exceed any fixed rational.
Together they describe both the fine structure and the unboundedness of
$\mathbb{Q}$.
\end{remark}

% ---------------------------------------------------------
\subsubsection{Farey Sequences}
% ---------------------------------------------------------

\begin{tcolorbox}[colback=defbox, colframe=defborder, arc=2pt,
  left=6pt, right=6pt, top=4pt, bottom=4pt,
  title={\small\textbf{Definition (Farey Sequence)}},
  fonttitle=\small\bfseries]
\begin{definition}[Farey sequence]
\label{def:farey-sequence}
For $n \in \mathbb{N}$, $n \geq 1$, the \emph{Farey sequence of order $n$},
denoted $F_n$, is the finite sequence of all rational numbers $a/b$ satisfying
\[
0 \leq \frac{a}{b} \leq 1, \qquad b \leq n, \qquad \gcd(a,b) = 1,
\]
arranged in strictly increasing order.
\end{definition}
\end{tcolorbox}

\begin{remark*}[Standard quantified statement]
\[
F_n = \left\{ \frac{a}{b} \in \mathbb{Q} \;\middle|\;
0 \leq a \leq b \leq n,\; \gcd(a,b)=1 \right\}
\text{ ordered increasingly.}
\]
\end{remark*}

\begin{remark*}[Interpretation]
$F_n$ is the canonical list of all fractions between $0$ and $1$ with
denominator at most $n$, each in reduced form.  As $n$ increases, new
fractions are inserted between existing neighbors.  The first few sequences are:
\begin{align*}
F_1 &= \Bigl(\tfrac{0}{1},\, \tfrac{1}{1}\Bigr), \\
F_2 &= \Bigl(\tfrac{0}{1},\, \tfrac{1}{2},\, \tfrac{1}{1}\Bigr), \\
F_3 &= \Bigl(\tfrac{0}{1},\, \tfrac{1}{3},\, \tfrac{1}{2},\,
              \tfrac{2}{3},\, \tfrac{1}{1}\Bigr), \\
F_4 &= \Bigl(\tfrac{0}{1},\, \tfrac{1}{4},\, \tfrac{1}{3},\,
              \tfrac{1}{2},\, \tfrac{2}{3},\, \tfrac{3}{4},\,
              \tfrac{1}{1}\Bigr).
\end{align*}
Each new fraction inserted between two Farey neighbors is their mediant,
and the neighbors of a newly inserted fraction always satisfy the Farey
determinant condition.
\end{remark*}

\begin{remark*}[Dependencies]
\hyperref[def:mediant]{Mediant},
\hyperref[def:rational-numbers]{Rational numbers},
\hyperref[thm:fraction-pair-equivalence]{reduced-form representatives}.
\end{remark*}

% ----- thm:farey-neighbor-theorem -----
\begin{tcolorbox}[colback=thmbox, colframe=thmborder, arc=2pt,
  left=6pt, right=6pt, top=4pt, bottom=4pt,
  title={\small\textbf{Theorem (Farey Neighbor Theorem)}},
  fonttitle=\small\bfseries]
\begin{theorem}[Farey Neighbor Theorem]
\label{thm:farey-neighbor-theorem}
Let $a/b < c/d$ be two consecutive fractions in the Farey sequence $F_n$,
with $\gcd(a,b) = \gcd(c,d) = 1$ and $b,d > 0$.  Then
\[
bc - ad = 1.
\]
\hyperref[prf:farey-neighbor-theorem]{Go to proof.}
\end{theorem}
\end{tcolorbox}

\begin{remark*}[Standard quantified statement]
\[
\forall \tfrac{a}{b}, \tfrac{c}{d} \in F_n,\quad
\tfrac{a}{b} \text{ and } \tfrac{c}{d} \text{ consecutive}
\;\Longrightarrow\;
bc - ad = 1.
\]
\end{remark*}

\begin{remark*}[Interpretation]
The determinant $bc - ad$ measures the ``arithmetic distance'' between two
reduced fractions.  For consecutive Farey neighbors this quantity is always
exactly $1$ — the smallest it can be for distinct reduced fractions with
positive denominators.  This is why Farey neighbors are best approximations
within their denominator range: no reduced fraction with smaller denominator
can fit between them.
\end{remark*}

\begin{remark*}[Negated quantified statement]
\[
\exists \tfrac{a}{b}, \tfrac{c}{d} \in F_n,\quad
\tfrac{a}{b} \text{ and } \tfrac{c}{d} \text{ consecutive}
\;\land\; bc - ad \neq 1.
\]
This negation is false: the theorem asserts the determinant is always $1$.
\end{remark*}

\begin{remark*}[Dependencies]
\begin{itemize}
  \item \hyperref[def:farey-sequence]{Farey Sequence}
  \item \hyperref[def:mediant]{Mediant}
\end{itemize}
\end{remark*}

% ----- thm:farey-gap -----
\begin{theorem}[Gap between Farey neighbors]
\label{thm:farey-gap}
Let $a/b < c/d$ be consecutive fractions in a Farey sequence.  Then
\[
\frac{c}{d} - \frac{a}{b} = \frac{1}{bd}.
\]
\hyperref[prf:farey-gap]{Go to proof.}
\end{theorem}

\begin{remark*}[Standard quantified statement]
\[
\forall \tfrac{a}{b}, \tfrac{c}{d} \text{ consecutive in some } F_n,\quad
\frac{c}{d} - \frac{a}{b} = \frac{1}{bd}.
\]
\end{remark*}

\begin{remark*}[Interpretation]
The gap between consecutive Farey neighbors is determined entirely by
their denominators.  Since $bd$ grows as $n$ increases, the gaps shrink —
this is the arithmetic explanation of why $\mathbb{Q}$ becomes arbitrarily
dense.  The identity also connects directly to Dirichlet approximation:
if $x$ lies between two Farey neighbors, it is within $1/(bd)$ of each,
giving a denominator-controlled approximation error.
\end{remark*}

\begin{remark*}[Dependencies]
\begin{itemize}
  \item \hyperref[thm:farey-neighbor-theorem]{Farey Neighbor Theorem}
\end{itemize}
\end{remark*}

% ----- thm:farey-mediant-property -----
\begin{theorem}[Mediant property of Farey neighbors]
\label{thm:farey-mediant-property}
If $a/b$ and $c/d$ are consecutive fractions in a Farey sequence, then
their mediant $(a+c)/(b+d)$ lies strictly between them:
\[
\frac{a}{b} < \frac{a+c}{b+d} < \frac{c}{d}.
\]
\hyperref[prf:farey-mediant-property]{Go to proof.}
\end{theorem}

\begin{remark*}[Standard quantified statement]
\[
\forall \tfrac{a}{b}, \tfrac{c}{d} \text{ consecutive Farey neighbors},\quad
\frac{a}{b} < \frac{a+c}{b+d} < \frac{c}{d}.
\]
\end{remark*}

\begin{remark*}[Interpretation]
This is the Mediant Inequality applied in the Farey context.  It guarantees
that the mediant insertion operation always produces a fraction that fits
strictly between consecutive neighbors, which is the mechanism by which
$F_n$ grows into $F_{n+1}$ (and by which the Stern--Brocot tree is generated).
\end{remark*}

\begin{remark*}[Dependencies]
\begin{itemize}
  \item \hyperref[thm:mediant-inequality]{Mediant Inequality}
  \item \hyperref[def:farey-sequence]{Farey Sequence}
\end{itemize}
\end{remark*}

% ----- lem:mediant-gap -----
\begin{lemma}[Mediant Gap Lemma]
\label{lem:mediant-gap}
If $a/b < c/d$ with $b, d > 0$, then the mediant $(a+c)/(b+d)$ lies
strictly between them, and the two subgaps satisfy
\[
\frac{a+c}{b+d} - \frac{a}{b} = \frac{1}{b(b+d)},
\qquad
\frac{c}{d} - \frac{a+c}{b+d} = \frac{1}{d(b+d)}.
\]
\hyperref[prf:mediant-gap]{Go to proof.}
\end{lemma}

\begin{remark*}[Standard quantified statement]
\[
\forall a,c \in \mathbb{Z},\; \forall b,d \in \mathbb{Z}_{>0},\quad
\frac{a}{b} < \frac{c}{d}
\;\Longrightarrow\;
\frac{a+c}{b+d} - \frac{a}{b} = \frac{1}{b(b+d)}
\;\land\;
\frac{c}{d} - \frac{a+c}{b+d} = \frac{1}{d(b+d)}.
\]
\end{remark*}

\begin{remark*}[Interpretation]
Inserting the mediant divides the interval $[a/b, c/d]$ into two
subintervals of lengths $1/(b(b+d))$ and $1/(d(b+d))$.  Both are
strictly smaller than the original gap $1/(bd)$ (since $b+d > b$ and
$b+d > d$).  This exact shrinkage formula is the engine behind the
Stern--Brocot limit theory: repeated mediant insertion forces interval
lengths to zero at a controlled rate determined entirely by the
denominators.
\end{remark*}

\begin{remark*}[Dependencies]
\begin{itemize}
  \item \hyperref[thm:mediant-inequality]{Mediant Inequality}
\end{itemize}
\end{remark*}

% ---------------------------------------------------------
\subsubsection{The Stern--Brocot Tree}
% ---------------------------------------------------------

\begin{theorem}[Stern--Brocot enumeration]
\label{thm:stern-brocot-enumeration}
Every positive rational number appears exactly once in reduced form in
the Stern--Brocot tree.
\hyperref[prf:stern-brocot-enumeration]{Go to proof.}
\end{theorem}

\begin{remark*}[Standard quantified statement]
\[
\forall q \in \mathbb{Q},\; q > 0,\;\;
\exists!\; \text{node in the Stern--Brocot tree whose reduced fraction equals } q.
\]
\end{remark*}

\begin{remark*}[Interpretation]
The Stern--Brocot tree is a complete, duplication-free enumeration of all
positive rationals.  Every positive rational is reachable by following a
unique finite left-right path from the root: go left if the target is less
than the current node, right if greater, and stop when equality holds.
The fact that this process always terminates (for rational targets) and
never revisits a number is a theorem — not a construction — and its proof
uses the Farey determinant condition together with a descent argument
analogous to the Euclidean algorithm.
\end{remark*}

\begin{remark*}[Dependencies]
\begin{itemize}
  \item \hyperref[def:mediant]{Mediant}
  \item \hyperref[thm:farey-neighbor-theorem]{Farey Neighbor Theorem}
  \item \hyperref[lem:mediant-gap]{Mediant Gap Lemma}
\end{itemize}
\end{remark*}

\begin{remark}[The Farey--Stern--Brocot connection]
The determinant condition $bc - ad = 1$ is the arithmetic shadow of the
Stern--Brocot tree.  Adjacent fractions in the tree always satisfy this
condition, and the mediant inserts the unique next fraction between them.
This is why the Farey gap formula $c/d - a/b = 1/(bd)$ and the Mediant
Gap Lemma control approximation errors so precisely: the entire arithmetic
structure of the tree is encoded in a single determinant identity.
\end{remark}
